\section{Defense-First Extension: Sources, Outcomes and Boundaries}\label{sec:r42-defense-first}
This purposive challenge set inspects four named seeds---CacheSolidarity/PrefixWall, PlanTwin, SICC and DP-MemView---plus two selected additions, OCELOT and AirGapAgent. Its purpose is to test scope interpretations against specific defense-oriented sources. The six family cards remain separate from the original 19-family/31-record C1--C5 cohort, which contains both attack-oriented and defense-oriented studies selected through empirical attack evaluation. Neither source frame estimates reporting prevalence across defense research.

\paragraph{Search history and incomplete coverage}
The retained materials record prior reconnaissance, nine logged queries and a cap of four seeds plus at most two additional relevant primary sources. Queries include exact identities, mitigation/isolation, padding/cache partitioning, agent privacy and a USENIX-focused check. The protocol is dated October~27, but has no independently verifiable preregistration timestamp. We report the cap as documentation of this bounded exercise, without claiming a prospectively registered stopping rule or inferring when it was chosen.

The search returns were not fully screened, and citation neighborhoods were not exhausted. The admission log retains unassessed candidates including AgentPrivArena and \emph{An AI Agent Execution Environment to Safeguard User Data} (GAAP); PriCon retains its prior abstract-only uncertainty. ISCA, MICRO, HPCA, ASPLOS and CHES/TCHES were not systematically covered by either search. SecEmb illustrates a consequence of the original eligibility rule: its GPT-2 quality/performance evidence remained adjacent because matched GPT-2 attack recovery was not established in the inspected scope \cite{umar2025secemb}. This is a selection-boundary example, not a completed hardware-defense search. The six cards therefore challenge particular interpretations without repairing omissions or selection bias.

\paragraph{A further hardware boundary check}
Two selected source checks illustrate evidence that the original empirical-attack selection need not admit. SecEmb argues protection for its specified embedding-access threat and reports GPT-2 quality and overhead, while its inspected attack demonstration uses a separate DLRM setting \cite[III--VI]{umar2025secemb}. Those results do not establish a matched GPT-2 attack-recovery evaluation. Palermo (\url{https://arxiv.org/html/2411.05400v1}) analyzes a GPT-2/OpenORCA workload in a modeled memory-system setting against an off-chip observer. Its binary response-time statistic concerns the specified stash/tree distinction; it is neither prompt recovery nor a general bound on all leakage. Its source describes simulated performance and controller area/power, while output quality and identical runs across these measurements were not established in the inspected scope. The old index includes SecEmb as adjacent; an exact-name and identifier search did not locate Palermo, which does not establish why it was absent or exclude an alias. These selected cases do not alter any archived cells or denominators and do not complete the extended search.

\paragraph{Source and outcome dictionary}
Each card records the exact inspected arXiv version, receiving permissions, intervention, attack/test condition, functional object, cost and transfer scope. R4 asks about protected-data attacker fitting; R5 requires clean-data and protected-data fits tested on the same defended population. A transcript-conditioned adaptive adversary is a separate capability. Functional quality applies to changed content or behavior; exact logical preservation supports only its stated object, leaving latency and task deadlines separate.

\paragraph{What the selected cards support}
All five quality-applicable cards locate a measured functional outcome; CacheSolidarity instead has narrow logical-preservation evidence. These positive examples do not form a homogeneous or randomly selected comparison with the original R8 count of 3/8 (4/8 accepting P). No population difference in reporting rates can be inferred from these counts. Four cards contain quantified empirical protection for their named receivers and two retain partial protection bindings. A functional score alone does not establish a same-configuration protection/function/cost package; the cards retain each binding separately. Semantic-disclosure, transcript-inference and auxiliary-channel outcomes also remain distinct. No complete R5 fit contrast was located in the inspected sections (four inapplicable, two unreported), while SICC and OCELOT explicitly evaluate broader adaptive attacks. This source-scoped R5 status must not be described as absence of adaptive evaluation.

\paragraph{Receiver-specific admission}
SICC isolates form-only length and structure inference, which warrants a distinct auxiliary-channel record candidate \cite[Table~3]{xu2026sicc}. DP-MemView's main endpoint is a semantic response transcript; its separate exposed internal-score diagnostic reports AUC 0.976 while the compliant response paths are fixed \cite[RQ2, after Table~4]{dpmemview2026}. That diagnostic must not be substituted for an attack against the compliant interface. PlanTwin's unrestricted schema re-identification is separate from budgeted quality; OCELOT's semantic trajectory guarantee is calibration-relative; AirGapAgent evaluates contextual disclosure rather than passive metadata \cite{yu2026plantwin,xie2026ocelot,bagdasarian2024airgap}.

\paragraph{Versions and runnable artifacts}
CacheSolidarity v1 is renamed PrefixWall in arXiv v2. Both inspected preprints retain a timing demonstration rather than a scored prompt-recovery estimate. The author-linked CacheSolidarity repository describes a richer Middleware artifact and a SafeKV reimplementation; its README expected results were not executed or substituted for the preprint \cite{pennas2026cachesolidarity}. PlanTwin's default heuristic description and SLM-backed timing table leave a configuration join to resolve \cite[5.1,5.4]{yu2026plantwin}. The source cards retain these uncertainties.

\onecolumn
\begingroup
\renewcommand{\arraystretch}{0.96}
\begin{longtable}{@{}p{.19\textwidth}p{.06\textwidth}p{.07\textwidth}p{.60\textwidth}@{}}
\caption{Complete source-scoped R1--R9 coding for the six defense-first cards (54 cells) and the separately scoped DP-MemView auxiliary-score diagnostic (nine cells). B = explicit binding, P = partial, U = unreported in inspected scope, ? = unverified, -- = inapplicable to the named outcome.}\label{tab:r42-defense-first-audit}\\
\toprule Card / source & Item & Code & Source locator / required scope\\\midrule\endfirsthead
\toprule Card / source & Item & Code & Source locator / required scope\\\midrule\endhead
DF01 CacheSolidarity / PrefixWall~\cite{pennas2026cachesolidarity} & R1 & B & 2.3;4\\
 & R2 & B & 2.2-2.3\\
 & R3 & B & 3.2-3.3;5.1\\
 & R4 & -- & 5.2.3:non-learned demonstration\\
 & R5 & -- & 5.2.3:no fit contrast\\
 & R6 & B & 5.1/Table2\\
 & R7 & P & 5.2.3/Fig9:timing trace, not scored recovery\\
 & R8 & -- & 3.2:narrow logical preservation\\
 & R9 & B & 5.2.4/Table3:component cost; retain object\\
\addlinespace[1pt]
DF02 PlanTwin~\cite{yu2026plantwin} & R1 & B & 3.2\\
 & R2 & B & 3.1-3.2\\
 & R3 & B & 4.1-4.4\\
 & R4 & -- & 5.6:field-overlap matcher\\
 & R5 & -- & 5.6:no fitted-attack contrast\\
 & R6 & B & 5.1\\
 & R7 & P & 5.2:exposure metric;5.6:unbudgeted inference\\
 & R8 & B & 5.2/Table5:planning quality only\\
 & R9 & P & 5.1 default heuristic versus5.4 SLM timing\\
\addlinespace[1pt]
DF03 SICC~\cite{xu2026sicc} & R1 & B & Task and Threat Model\\
 & R2 & B & Task and Threat Model\\
 & R3 & B & Operational Compiler\\
 & R4 & P & Setup:train/test stated; defended-fit allocation unclear\\
 & R5 & U & Setup/Adaptive Attacks:no matched clean/protected fit contrast\\
 & R6 & B & Setup;Tables7-9\\
 & R7 & B & Tables3,5,6,8\\
 & R8 & B & Tables7-9:condition-matched completion\\
 & R9 & P & Operational cost:operation description only\\
\addlinespace[1pt]
DF04 DP-MemView~\cite{dpmemview2026} & R1 & B & Problem Definition\\
 & R2 & B & Problem Definition/interface contract\\
 & R3 & B & DP-MemView;Table1\\
 & R4 & P & C.5:fit/calibration split; per-interface fitting unclear\\
 & R5 & U & Experiments/C.5:no clean/protected fit contrast\\
 & R6 & B & C.1-C.3/Table5\\
 & R7 & B & Table2;C.5:empirical transcript AUC\\
 & R8 & B & Table2;C.5:matched model-judged quality\\
 & R9 & U & Experiments/C.5:no operational cost measurement\\
\addlinespace[1pt]
DF05 OCELOT~\cite{xie2026ocelot} & R1 & B & III-C-D\\
 & R2 & B & III-B-C\\
 & R3 & B & IV;Algorithm1\\
 & R4 & -- & V-A:fixed LLM proxy pools;defense calibration separate\\
 & R5 & -- & V-A,V-E:no clean/protected fitted-attacker comparison\\
 & R6 & B & V-A;AppendixD\\
 & R7 & B & V-B/TableII;V-E/TableV\\
 & R8 & B & V-B/TableII;V-F/TableVI\\
 & R9 & B & V-F/TableVII\\
\addlinespace[1pt]
DF06 AirGapAgent~\cite{bagdasarian2024airgap} & R1 & B & 2.2-2.3\\
 & R2 & B & 2.3\\
 & R3 & B & 4.1-4.2\\
 & R4 & -- & 6.1:prompted models,not learned attacker fitting\\
 & R5 & -- & 5.3;6:no clean/protected fitted-attacker comparison\\
 & R6 & B & 6.1\\
 & R7 & B & 5.3;Tables3-6:explicit-disclosure protection\\
 & R8 & B & Tables3-5:matched answer utility\\
 & R9 & U & 6:cost not located\\
\addlinespace[1pt]
DF04-score~\cite{dpmemview2026} & R1 & B & RQ2 Interface Safeguards, after Table4;B.5\\
 & R2 & B & RQ2 Interface Safeguards, after Table4;B.5\\
 & R3 & B & RQ2 Interface Safeguards, after Table4;B.5\\
 & R4 & -- & RQ2 Interface Safeguards, after Table4;B.5\\
 & R5 & -- & RQ2 Interface Safeguards, after Table4;B.5\\
 & R6 & P & RQ2 Interface Safeguards, after Table4;B.5; exact diagnostic denominator is partial\\
 & R7 & B & RQ2 Interface Safeguards, after Table4;B.5; exposed top-score AUC 0.976\\
 & R8 & -- & RQ2 Interface Safeguards, after Table4;B.5; compliant response paths fixed; narrow preservation only\\
 & R9 & U & RQ2 Interface Safeguards, after Table4;B.5\\
\bottomrule\end{longtable}

\endgroup
\twocolumn

The artifact directory \path{data/r42-defense-first/} retains the protocol, search returns, admission decisions, source-read scopes, cards and deterministic projection. This extension records source-reported results; new model and defense executions are documented separately.
